% ===========================================================================
% mappe/blaetter.tex — die vier Blätter der Charaktermappe
%
% Jedes Blatt ist eine mappenblatt-Umgebung und damit eine Box von genau
% 210 x 297 mm. Wo sie landet, entscheidet mappe.tex: eine pro A4-Seite
% (digital) oder zwei nebeneinander auf A3 quer (Druck).
%
% Alle Maße sind am HvS-Charakterbogen abgelesen — die Textzeilen mit
% pdftotext -bbox, Titel, Figur und Kasten über eine Pixelanalyse des
% 300-dpi-Rasters:
%
%   Titelblatt    Titel-Lockup  y 19,2..63,9   x 50,8..159,0 (Mitte 105,0)
%                 Figur         y 75..244      x 71,5..156,2 (85 breit)
%                 Fußtext       y 255,9 ff.    x 20..192   Zeilenabstand 5,93
%   Rückseite     Überschrift   y 28,5         mittig
%                 Fließtext     y 39,1 ff.     x 30,6..179,4  Abstand 6,77
%                 Fußkasten     Titel y 244,8, Text darunter
%
% Der Titel ist EIN Lockup, kein Name plus separat gesetzte Profession:
% beim Vorbild sind die drei Zeilen ein Block von 19,2 bis 63,9 mm. Deshalb
% steht die Profession als dritte Zeile im selben Knoten. Der Preis dafür
% ist die \fontsize-Klammer mitten im Text; der Gewinn ist, dass ein
% zweizeiliger Name die Profession von selbst nach unten schiebt.
% ===========================================================================

% --- Blatt 1: Titel --------------------------------------------------------
\newcommand{\blatttitel}{%
  \begin{mappenblatt}
    \flechtrahmen
    % Kopfzier statt des Verlagslogos — siehe mechanik-mappe.tex
    \kopfzier{105}{14}{86}{4}
    \konturzeile{105}{21}{0.30}%
      {\auszeichnung\fontsize{\mappenamengroesse}{\mappenamenzeile}\selectfont}%
      {\wertoder{mappe_name}{s1_name}%
       \wertdaP{mappe_profession}{%
         \\[2.5mm]{\fontsize{\mappeprofgroesse}{\mappeprofgroesse}\selectfont
                   \wertvon{mappe_profession}}%
       }{}}%
    % Ganzkörperbild, auf der Standlinie
    \mappebild{mappe_bild}{s1_bild}{62.5}{75}{85}{169}
    % Fußkasten mit der Kurzempfehlung
    \wertdaP{mappe_empfehlung}{%
      \pergamentkasten{14}{248}{182}{%
        \lesetext\fontsize{11.5}{15.5}\selectfont\centering
        \wertvon{mappe_empfehlung}}%
    }{}%
  \end{mappenblatt}}

% --- Blatt 2: Rückseite ----------------------------------------------------
\newcommand{\blattrueckseite}{%
  \begin{mappenblatt}
    % Das Heldenbild als blasses Wasserzeichen hinter dem Text, wie beim
    % Vorbild: dieselbe Stelle, nur heruntergenommen.
    \begin{scope}[opacity=0.13]
      \mappebild{mappe_bild}{s1_bild}{62.5}{75}{85}{169}
    \end{scope}
    \flechtrahmen
    \node[anchor=north,text=schrift,font=\auszeichnung\fontsize{19}{23}\selectfont]
      at (105,28.5) {\mappehintergrundtitel};
    \mappetextblock{30.6}{38}{148.9}%
      {\lesetext\fontsize{13}{19.2}\selectfont}{justify}%
      {\wertvon{mappe_hintergrund}}%
    % Fußkasten, frei betextbar
    \wertdaP{mappe_fusstext}{%
      \pergamentkasten{14}{243}{182}{%
        \centering
        \wertdaP{mappe_fusstitel}{%
          {\auszeichnung\fontsize{13}{17}\selectfont
           \wertvon{mappe_fusstitel}\par\vspace{1.4mm}}%
        }{}%
        \lesetext\fontsize{11.5}{15.5}\selectfont
        \wertvon{mappe_fusstext}}%
    }{}%
  \end{mappenblatt}}

% --- Blatt 3: Innenseite links ---------------------------------------------
% Trägt den Pflichttext der Vereinbarung über Gemeinschaftsinhalte. Der muss
% laut Baukasten ("Lies mich zuerst", Seite 2: "Der folgende Text muss stets
% im Werk enthalten sein") wortgleich im Werk stehen — deshalb steht er nicht
% in der Heldendatei. Nur die Copyright-Zeile ist variabel.
%
% Der Wortlaut selbst steht in ../pflichttext.tex, gemeinsam mit dem
% Impressum der Abenteuerklasse. Hier stand eine ältere Fassung mit einem
% anderen Rechteinhaber ("von Ulisses Medien und Spiele Distribution GmbH in
% Deutschland, den U.S.A. und anderen Ländern") und einer anderen
% Markenliste, während dsa5latex.cls die aktuelle führte. Welche gilt,
% entscheidet sich jetzt an einer Stelle.
%
% Die Absatztrennung bleibt hier bei der Vorgabe \par: den Abstand macht
% \parskip, siehe unten.
\newcommand{\blattinnenlinks}{%
  \begin{mappenblatt}
    \flechtrahmen
    \node[anchor=north,text=schrift,font=\auszeichnung\fontsize{15}{19}\selectfont]
      at (105,28.5) {RECHTLICHES};
    % Absätze durch Leerzeilen, nicht durch \par\vspace: nach einem \vspace
    % setzt TikZ den folgenden Absatz im Blocksatz statt zentriert, und dann
    % ist ein Absatz von dreien anders ausgerichtet als die übrigen.
    \mappetextblock{30.6}{42}{148.9}%
      {\lesetext\fontsize{10.5}{14.5}\selectfont\setlength{\parskip}{4mm}}%
      {center}{\dsaPflichttext{\mappecopyright}}%
    \wertdaP{mappe_innen_links_text}{%
      \pergamentkasten{14}{205}{182}{%
        \centering
        \wertdaP{mappe_innen_links_titel}{%
          {\auszeichnung\fontsize{13}{17}\selectfont
           \wertvon{mappe_innen_links_titel}\par\vspace{1.4mm}}%
        }{}%
        \lesetext\fontsize{11.5}{15.5}\selectfont
        \wertvon{mappe_innen_links_text}}%
    }{}%
  \end{mappenblatt}}

% --- Blatt 4: Innenseite rechts --------------------------------------------
% Frei betextbar. Ohne Text: Notizlinien, damit die Seite am Spieltisch
% etwas taugt und nicht bloß leer ist.
\newcommand{\blattinnenrechts}{%
  \begin{mappenblatt}
    \flechtrahmen
    \node[anchor=north,text=schrift,font=\auszeichnung\fontsize{15}{19}\selectfont]
      at (105,28.5) {%
      \wertdaP{mappe_innen_rechts_titel}{\wertvon{mappe_innen_rechts_titel}}{NOTIZEN}};
    \wertdaP{mappe_innen_rechts_text}{%
      \mappetextblock{30.6}{42}{148.9}%
        {\lesetext\fontsize{12}{17.5}\selectfont}{justify}%
        {\wertvon{mappe_innen_rechts_text}}%
    }{%
      \foreach \i in {0,...,24} {%
        \draw[bandkontur,opacity=0.26,line width=0.25mm]
          (26,{46+\i*9}) -- (184,{46+\i*9});}%
    }%
  \end{mappenblatt}}
